\documentclass[11pt,a4paper]{article}

\usepackage[italian]{babel}
\usepackage[T1]{fontenc}
\usepackage[utf8]{inputenc}
\usepackage{amsmath,amssymb,amsthm}
\usepackage{mathtools}
\usepackage{booktabs}
\usepackage{microtype}
\usepackage{geometry}
\usepackage[hidelinks]{hyperref}

\geometry{margin=3cm}

\newtheorem{theorem}{Teorema}[section]
\newtheorem{proposition}[theorem]{Proposizione}
\newtheorem{lemma}[theorem]{Lemma}
\newtheorem{corollary}[theorem]{Corollario}
\theoremstyle{definition}
\newtheorem{definition}[theorem]{Definizione}
\theoremstyle{remark}
\newtheorem{remark}[theorem]{Osservazione}
\newtheorem{example}[theorem]{Esempio}

\newcommand{\Db}{D_b}
\newcommand{\MSC}{\mathcal M}
\newcommand{\val}{\operatorname{val}}
\newcommand{\ord}{\operatorname{ord}}
\newcommand{\fix}{\operatorname{fix}}
\newcommand{\mcm}{\operatorname{mcm}}

\title{Permutazioni digitalmente separabili nei giochi di carte su
	\texorpdfstring{\(b^m\)}{b alla m} posizioni\\
	\large Prodotti di Kronecker e ricostruzione da fibre digitali}
\author{Maurizio Berlanda\\
	\small\url{https://github.com/maberlanda/27-card-tensor-structure}}
\date{}

\begin{document}
	
	\maketitle
	
	\begin{abstract}
		Si considera una famiglia di procedure su \(N=b^m\) posizioni, ottenute
		ripetendo per \(m\) stadi la distribuzione ciclica in \(b\) mazzetti e la
		loro raccolta secondo una permutazione assegnata. La scrittura in base \(b\)
		identifica le posizioni con le parole di \(D_b^m\), dove
		\(D_b=\{0,\dots,b-1\}\), e descrive il ciclo completo mediante un'azione
		separata sulle cifre. La trasformazione globale possiede la fattorizzazione
		matriciale
		\[
		P_{\boldsymbol{\rho}}
		=
		P_{m-1}\otimes\cdots\otimes P_0.
		\]
		Si dimostra una caratterizzazione intrinseca delle permutazioni così
		ottenute: esse sono esattamente le permutazioni che preservano, per ogni
		livello, la partizione delle posizioni determinata dal valore della cifra
		corrispondente. La caratterizzazione fornisce l'unicità dei fattori locali e
		realizza un'immersione \(S_b^m\hookrightarrow S_{b^m}\). Si introducono poi
		statistiche aggregate sulle fibre digitali e si prova che, all'interno di
		questa classe, esse determinano integralmente le permutazioni locali. Il caso
		\(b=m=3\) viene sviluppato esplicitamente e collegato al gioco delle
		ventisette carte. Un'appendice confronta infine la struttura omogenea
		\(b^m\) con il caso rettangolare \(N=ij\) e con la contrazione informativa
		del gioco delle ventuno carte.
	\end{abstract}
	
	\section{Introduzione}
	
	Il gioco delle ventisette carte appartiene alla famiglia dei procedimenti
	classici nei quali una distribuzione ripetuta traduce informazioni operative
	in cifre di una scrittura posizionale. L'identità
	\[
	27=3^3
	\]
	permette di distribuire il mazzo in tre mazzetti e di ripetere il gesto per
	tre stadi; le posizioni occupate, nelle raccolte successive, dai mazzetti che
	contengono una carta determinano progressivamente le cifre ternarie della sua
	posizione finale. Il legame con la notazione posizionale e con gli algoritmi
	di ordinamento per cifre è classico nella letteratura sul trucco di Gergonne
	\cite{bolker}.
	
	Il modello qui studiato isola una versione non adattiva del procedimento. Si
	fissano interi
	\[
	b\ge 2,
	\qquad
	m\ge 2,
	\]
	e un mazzo di \(N=b^m\) carte. In ciascuno dei \(m\) stadi il mazzo viene
	distribuito ciclicamente in \(b\) mazzetti, conservando l'ordine interno, e i
	mazzetti vengono raccolti secondo una permutazione di \(S_b\) stabilita per
	quello stadio. Non sono ammessi rovesciamenti del mazzo né scelte adattive
	dipendenti dalla posizione di una carta. La successione delle permutazioni di
	raccolta forma un codice locale
	\[
	\boldsymbol{\rho}
	=
	(\rho_{m-1},\dots,\rho_0)
	\in S_b^m,
	\]
	che determina una permutazione globale delle \(b^m\) posizioni.
	
	La domanda centrale non riguarda soltanto il calcolo di questa permutazione.
	Occorre riconoscere, fra tutti gli elementi di \(S_{b^m}\), quelli che possono
	nascere dal procedimento, individuare la fattorizzazione locale nascosta nella
	loro azione e stabilire quali informazioni aggregate siano sufficienti a
	ricostruirla. La scrittura posizionale, il prodotto tensoriale e l'azione del
	gruppo diretto \(S_b^m\) convergono così nello stesso problema.
	
	Il ciclo fisico viene dapprima descritto mediante una rotazione dei fattori
	tensoriali; dopo \(m\) stadi la rotazione si chiude e la matrice globale si
	separa in un prodotto di Kronecker. Il risultato conduce a una
	caratterizzazione intrinseca: una permutazione globale appartiene alla classe
	considerata esattamente quando preserva, livello per livello, le partizioni
	digitali delle posizioni. Le permutazioni locali sono allora determinate in
	modo unico. Su questa classe, le somme delle immagini lungo le fibre digitali
	costituiscono inoltre un insieme di statistiche aggregate sufficiente alla
	ricostruzione completa del codice.
	
	Il formalismo dei prodotti di Kronecker e delle matrici di commutazione è
	standard; per le identità fondamentali si rinvia a Henderson e Searle
	\cite{henderson-searle}. La teoria dei perfect shuffles offre un contesto più
	ampio per lo studio algebrico delle permutazioni di mazzi \cite{diaconis-graham-kantor},
	mentre analisi dinamiche di giochi di carte affini compaiono, fra gli altri,
	in Champanerkar e Jani \cite{champanerkar-jani} e Quintero \cite{quintero}.
	
	\section{Modello posizionale e procedura ordinaria}
	
	Sia
	\[
	\Db=\{0,1,\dots,b-1\}.
	\]
	Le posizioni del mazzo sono numerate da \(0\) a \(b^m-1\). Ogni posizione
	possiede una rappresentazione unica
	\[
	x=(a_{m-1}a_{m-2}\cdots a_1a_0)_b,
	\qquad
	a_i\in\Db,
	\]
	ossia
	\[
	x=a_{m-1}b^{m-1}+a_{m-2}b^{m-2}+\cdots+a_1b+a_0.
	\]
	Per
	\[
	\alpha=(a_{m-1},\dots,a_0)\in\Db^m
	\]
	si pone
	\[
	\val(\alpha)
	=
	\sum_{i=0}^{m-1}a_i b^i.
	\]
	La funzione \(\val\) è una biiezione fra \(\Db^m\) e
	\(\{0,\dots,b^m-1\}\).
	
	Durante un singolo stadio, la posizione
	\[
	x=bq+r,
	\qquad
	r=x\bmod b,
	\qquad
	q=\left\lfloor\frac{x}{b}\right\rfloor
	\]
	viene distribuita nel mazzetto identificato dal resto \(r\). Se una
	permutazione \(\rho\in S_b\) colloca quel mazzetto nella posizione di blocco
	\(\rho(r)\), la nuova posizione è
	\[
	x'
	=
	\rho(r)b^{m-1}
	+
	\left\lfloor\frac{x}{b}\right\rfloor.
	\]
	La cifra meno significativa viene così estratta dall'indirizzo e reinserita
	come cifra più significativa, dopo l'azione della permutazione di raccolta.
	
	\begin{proposition}[Ricorrenza posizionale]
		Siano \(\rho_0,\dots,\rho_{k-1}\in S_b\) le permutazioni usate nei primi
		\(k\) stadi, con \(1\le k\le m\), e sia
		\[
		x_0=(a_{m-1}\cdots a_0)_b.
		\]
		Allora
		\[
		\begin{aligned}
			x_k
			={}&
			\rho_{k-1}(a_{k-1})b^{m-1}
			+\rho_{k-2}(a_{k-2})b^{m-2}
			+\cdots
			+\rho_0(a_0)b^{m-k}
			\\
			&+
			\left\lfloor\frac{x_0}{b^k}\right\rfloor.
		\end{aligned}
		\]
		In particolare, dopo \(m\) stadi,
		\[
		x_m
		=
		\sum_{i=0}^{m-1}\rho_i(a_i)b^i.
		\]
	\end{proposition}
	
	\begin{proof}
		Per \(k=1\) la formula coincide con la regola del singolo stadio. Se essa
		vale per un certo \(k<m\), la cifra successivamente estratta dal termine
		\(\lfloor x_0/b^k\rfloor\) è \(a_k\), mentre ciascuno dei termini già
		prodotti è divisibile per \(b\). La relazione
		\[
		x_{k+1}
		=
		\rho_k(a_k)b^{m-1}
		+
		\left\lfloor\frac{x_k}{b}\right\rfloor
		\]
		fornisce quindi la formula con \(k+1\) al posto di \(k\). Per \(k=m\), il
		termine \(\lfloor x_0/b^m\rfloor\) è nullo e le cifre trasformate occupano i
		livelli originari.
	\end{proof}
	
	Nel caso \(b=m=3\), la formula descrive il ciclo completo del gioco delle
	ventisette carte. Il suo significato tensoriale permette di distinguere con
	precisione l'operatore di uno stadio dalla trasformazione finale.
	
	\section{Dinamica tensoriale del ciclo completo}
	
	Sia
	\[
	E_b=\mathbb R^b
	\]
	con base canonica \(e_0,\dots,e_{b-1}\). All'indirizzo
	\(\alpha=(a_{m-1},\dots,a_0)\) si associa il tensore
	\[
	e_\alpha
	=
	e_{a_{m-1}}\otimes\cdots\otimes e_{a_0}
	\in E_b^{\otimes m}.
	\]
	La distribuzione ciclica è rappresentata dall'operatore
	\[
	\MSC_m:E_b^{\otimes m}\longrightarrow E_b^{\otimes m}
	\]
	definito sui tensori di base da
	\[
	\MSC_m
	\left(
	e_{a_{m-1}}\otimes\cdots\otimes e_{a_0}
	\right)
	=
	e_{a_0}\otimes e_{a_{m-1}}\otimes\cdots\otimes e_{a_1}.
	\]
	Esso realizza una rotazione ciclica dei fattori e soddisfa
	\[
	\MSC_m^m=I.
	\]
	
	Sia \(P_h\) la matrice di permutazione associata a \(\rho_h\), secondo la
	convenzione
	\[
	P_he_r=e_{\rho_h(r)}.
	\]
	L'operatore dello stadio \(h\), con \(h=0,\dots,m-1\), è
	\[
	T_h
	=
	(P_h\otimes I_b^{\otimes(m-1)})\circ\MSC_m.
	\]
	La rotazione porta la cifra da elaborare nel primo fattore; la raccolta
	applica poi \(P_h\) a quel fattore. La separazione delle azioni locali emerge
	soltanto dopo il ciclo completo.
	
	\begin{theorem}[Chiusura tensoriale del ciclo]
		Con le convenzioni precedenti,
		\[
		T_{m-1}\circ\cdots\circ T_1\circ T_0
		=
		P_{m-1}\otimes P_{m-2}\otimes\cdots\otimes P_0.
		\]
	\end{theorem}
	
	\begin{proof}
		Applicando successivamente gli operatori a
		\[
		e_{a_{m-1}}\otimes\cdots\otimes e_{a_0},
		\]
		dopo il primo stadio si ottiene
		\[
		e_{\rho_0(a_0)}\otimes e_{a_{m-1}}\otimes\cdots\otimes e_{a_1},
		\]
		e dopo il secondo
		\[
		e_{\rho_1(a_1)}\otimes e_{\rho_0(a_0)}
		\otimes e_{a_{m-1}}\otimes\cdots\otimes e_{a_2}.
		\]
		Proseguendo per \(m\) stadi, ciascuna cifra viene esposta una volta
		all'azione della permutazione corrispondente e i fattori ritornano
		nell'ordine iniziale. Il risultato è
		\[
		e_{\rho_{m-1}(a_{m-1})}
		\otimes\cdots\otimes
		e_{\rho_0(a_0)},
		\]
		che coincide con l'azione di
		\(P_{m-1}\otimes\cdots\otimes P_0\). I tensori elementari formano una base,
		perciò i due operatori sono uguali.
	\end{proof}
	
	La relazione \(\MSC_m^m=I\) non autorizza una cancellazione formale delle
	rotazioni nella composizione, poiché fra due rotazioni successive compaiono
	le trasformazioni locali. Il teorema descrive invece il trasporto di ogni
	fattore fino al livello sul quale deve agire la raccolta e la successiva
	chiusura dell'intero ciclo.
	
	\section{Permutazioni digitalmente separabili}
	
	Il codice locale
	\[
	\boldsymbol{\rho}
	=
	(\rho_{m-1},\dots,\rho_0)
	\in S_b^m
	\]
	definisce la trasformazione globale
	\[
	\pi_{\boldsymbol{\rho}}
	\bigl(\val(a_{m-1},\dots,a_0)\bigr)
	=
	\val\bigl(
	\rho_{m-1}(a_{m-1}),\dots,\rho_0(a_0)
	\bigr).
	\]
	La matrice di permutazione corrispondente, nella base ordinata dai tensori
	\(e_\alpha\), è
	\[
	P_{\boldsymbol{\rho}}
	=
	P_{m-1}\otimes\cdots\otimes P_0.
	\]
	
	Per ogni livello \(i\in\{0,\dots,m-1\}\), sia
	\[
	d_i\bigl(\val(a_{m-1},\dots,a_0)\bigr)=a_i
	\]
	la proiezione sulla cifra di livello \(i\). Per \(t\in\Db\) si definisce la
	fibra digitale
	\[
	F_{i,t}
	=
	\{x\in\{0,\dots,b^m-1\}:d_i(x)=t\},
	\]
	e la partizione
	\[
	\mathcal F_i
	=
	\{F_{i,t}:t\in\Db\}.
	\]
	Ogni fibra ha cardinalità \(b^{m-1}\).
	
	\begin{definition}
		Una permutazione \(\pi\in S_{b^m}\) è \emph{digitalmente separabile} se
		esistono permutazioni \(\rho_0,\dots,\rho_{m-1}\in S_b\) tali che
		\[
		d_i(\pi(x))
		=
		\rho_i(d_i(x))
		\]
		per ogni posizione \(x\) e ogni livello \(i\).
	\end{definition}
	
	La definizione può essere riconosciuta direttamente nell'azione di \(\pi\),
	senza conoscere in anticipo il codice che l'ha prodotta.
	
	\begin{theorem}[Caratterizzazione intrinseca e fattorizzazione]
		Per una permutazione \(\pi\in S_{b^m}\), le seguenti proprietà sono
		equivalenti.
		\begin{enumerate}
			\item Esiste un codice unico
			\[
			\boldsymbol{\rho}
			=(\rho_{m-1},\dots,\rho_0)\in S_b^m
			\]
			tale che \(\pi=\pi_{\boldsymbol{\rho}}\).
			
			\item Per ogni \(i\), la funzione \(d_i\circ\pi\) è costante su ciascuna
			fibra \(F_{i,t}\).
			
			\item Per ogni \(i\), la permutazione \(\pi\) induce una permutazione dei
			blocchi della partizione \(\mathcal F_i\).
			
			\item La matrice di permutazione di \(\pi\) ammette una fattorizzazione
			\[
			P_{m-1}\otimes\cdots\otimes P_0,
			\]
			dove ciascun \(P_i\) è una matrice di permutazione \(b\times b\).
		\end{enumerate}
		Le permutazioni \(\rho_i\) e le matrici \(P_i\) sono determinate in modo
		unico.
	\end{theorem}
	
	\begin{proof}
		La prima proprietà implica la seconda perché, per
		\(x=\val(a_{m-1},\dots,a_0)\), vale
		\[
		d_i(\pi(x))=\rho_i(a_i).
		\]
		Il membro destro dipende soltanto dal valore \(a_i=d_i(x)\), perciò è
		costante su \(F_{i,t}\).
		
		Supposta la seconda proprietà, per ogni \(i\) e \(t\) si definisce
		\[
		\rho_i(t)
		=
		d_i(\pi(x)),
		\qquad x\in F_{i,t}.
		\]
		La definizione non dipende dalla scelta di \(x\). Se due valori distinti
		\(t,t'\) avessero la stessa immagine \(u\), l'iniettività di \(\pi\) dovrebbe
		mandare i due insiemi disgiunti \(F_{i,t}\) e \(F_{i,t'}\), ciascuno di
		cardinalità \(b^{m-1}\), dentro la sola fibra \(F_{i,u}\), che possiede la
		stessa cardinalità di uno di essi. Ciò è impossibile. La funzione \(\rho_i\)
		è quindi iniettiva e, poiché \(\Db\) è finito, appartiene a \(S_b\).
		
		Per ogni posizione \(x\), le cifre di \(\pi(x)\) sono allora
		\[
		\rho_{m-1}(d_{m-1}(x)),\dots,\rho_0(d_0(x)).
		\]
		Esse determinano un'unica posizione, e ne segue
		\(\pi=\pi_{\boldsymbol{\rho}}\). La stessa costruzione mostra l'unicità di
		ciascuna \(\rho_i\).
		
		La seconda e la terza proprietà sono equivalenti: la costanza di
		\(d_i\circ\pi\) su \(F_{i,t}\) equivale all'inclusione dell'immagine della
		fibra in una singola fibra finale; la biiettività di \(\pi\) e l'uguaglianza
		delle cardinalità trasformano l'inclusione in uguaglianza.
		
		Infine, la prima proprietà implica la quarta per il teorema di chiusura
		tensoriale. Viceversa, una matrice
		\(P_{m-1}\otimes\cdots\otimes P_0\) agisce sui tensori di base trasformando
		separatamente i rispettivi indici, e induce quindi
		\(\pi_{\boldsymbol{\rho}}\). L'unicità delle permutazioni locali comporta
		quella dei fattori matriciali all'interno della classe delle matrici di
		permutazione.
	\end{proof}
	
	\begin{corollary}
		L'applicazione
		\[
		\Phi:S_b^m\longrightarrow S_{b^m},
		\qquad
		\boldsymbol{\rho}\longmapsto\pi_{\boldsymbol{\rho}},
		\]
		è un omomorfismo iniettivo. La sua immagine coincide con il gruppo delle
		permutazioni digitalmente separabili ed è isomorfa a \(S_b^m\).
	\end{corollary}
	
	\begin{proof}
		La composizione di due codici avviene componente per componente, perciò
		\[
		\pi_{\boldsymbol{\rho}\boldsymbol{\sigma}}
		=
		\pi_{\boldsymbol{\rho}}\circ\pi_{\boldsymbol{\sigma}}.
		\]
		L'unicità del codice associato a una trasformazione rende il nucleo banale.
		La descrizione dell'immagine segue dal teorema.
	\end{proof}
	
	Nel caso \(b=m=3\), si ottiene
	\[
	S_3^3\hookrightarrow S_{27}.
	\]
	L'immagine contiene \(6^3=216\) permutazioni, una sottoclasse fortemente
	vincolata dei \(27!\) elementi di \(S_{27}\).
	
	\section{Ricostruzione da statistiche aggregate sulle fibre}
	
	La caratterizzazione precedente ricostruisce le componenti locali leggendo le
	cifre delle immagini. Un secondo problema riguarda la quantità di
	informazione necessaria quando, al posto dell'intera permutazione, siano
	disponibili dati aggregati sulle fibre digitali.
	
	Per una permutazione digitalmente separabile
	\(\pi=\pi_{\boldsymbol{\rho}}\), si definisce
	\[
	\Sigma_{i,t}(\pi)
	=
	\sum_{x\in F_{i,t}}\pi(x).
	\]
	La somma raccoglie le immagini di tutte le posizioni che condividono la cifra
	\(t\) al livello \(i\).
	
	\begin{theorem}[Formula delle statistiche di fibra]
		Per ogni \(i\in\{0,\dots,m-1\}\) e \(t\in\Db\),
		\[
		\Sigma_{i,t}(\pi)
		=
		C_i+b^{m-1+i}\rho_i(t),
		\]
		dove
		\[
		C_i
		=
		b^{m-2}\frac{b(b-1)}2
		\sum_{\substack{0\le j\le m-1\\j\ne i}}b^j.
		\]
	\end{theorem}
	
	\begin{proof}
		Se \(x=\val(a_{m-1},\dots,a_0)\), allora
		\[
		\pi(x)
		=
		\sum_{j=0}^{m-1}b^j\rho_j(a_j).
		\]
		Sulla fibra \(F_{i,t}\), la cifra \(a_i\) è fissata. Il contributo del
		livello \(i\) alla somma è quindi
		\[
		\sum_{x\in F_{i,t}}b^i\rho_i(t)
		=
		b^{m-1+i}\rho_i(t).
		\]
		Per ogni \(j\ne i\), la cifra \(a_j\) assume ciascun valore di \(\Db\)
		esattamente \(b^{m-2}\) volte. Poiché \(\rho_j\) è una permutazione,
		\[
		\sum_{u=0}^{b-1}\rho_j(u)
		=
		\sum_{u=0}^{b-1}u
		=
		\frac{b(b-1)}2.
		\]
		Sommando i contributi dei livelli \(j\ne i\) si ottiene \(C_i\), che non
		dipende da \(t\) né dalle permutazioni locali diverse da \(\rho_i\).
	\end{proof}
	
	\begin{corollary}[Sufficienza delle statistiche aggregate]
		All'interno del gruppo delle permutazioni digitalmente separabili, la
		famiglia completa delle quantità \(\Sigma_{i,t}\) determina univocamente la
		permutazione globale. Le componenti locali sono ricostruite mediante
		\[
		\rho_i(t)
		=
		\frac{\Sigma_{i,t}(\pi)-C_i}{b^{m-1+i}}.
		\]
		In particolare, la mappa
		\[
		\pi
		\longmapsto
		\bigl(\Sigma_{i,t}(\pi)\bigr)_{
			0\le i\le m-1,\;t\in\Db}
		\]
		è iniettiva sulla classe digitalmente separabile.
	\end{corollary}
	
	\begin{proof}
		La formula ricostruisce tutti i valori di tutte le permutazioni
		\(\rho_i\). Il teorema di caratterizzazione determina quindi
		\(\pi=\pi_{\boldsymbol{\rho}}\) in modo unico.
	\end{proof}
	
	\begin{remark}
		Se l'intera permutazione \(\pi\) è conosciuta punto per punto, le componenti
		\(\rho_i\) possono essere lette direttamente dalle cifre di immagini
		opportunamente scelte. Le somme sulle fibre non costituiscono dunque l'unico
		metodo di ricostruzione; mostrano piuttosto che un insieme di statistiche
		aggregate conserva tutta l'informazione locale entro la classe separabile.
		La sola compatibilità numerica delle somme con la formula non viene qui
		assunta come criterio di riconoscimento per una permutazione arbitraria di
		\(S_{b^m}\).
	\end{remark}
	
	\section{Il caso delle ventisette carte}
	
	Per \(b=m=3\), le posizioni sono numerate da \(0\) a \(26\) e ogni posizione
	ha la forma
	\[
	x=(a_2a_1a_0)_3.
	\]
	Le fibre digitali contengono nove posizioni. La formula generale diventa
	\[
	\Sigma_{i,t}(\pi)
	=
	C_i+3^{2+i}\rho_i(t),
	\]
	con
	\[
	C_0=108,
	\qquad
	C_1=90,
	\qquad
	C_2=36.
	\]
	I valori possibili sono
	\[
	\begin{array}{c@{\qquad}ccc}
		\toprule
		i & \rho_i(t)=0 & \rho_i(t)=1 & \rho_i(t)=2\\
		\midrule
		0 & 108 & 117 & 126\\
		1 &  90 & 117 & 144\\
		2 &  36 & 117 & 198\\
		\bottomrule
	\end{array}
	\]
	Le differenze consecutive sono \(9\), \(27\) e \(81\), in accordo con i
	coefficienti \(3^{2+i}\). Il valore comune \(117\), ottenuto quando
	\(\rho_i(t)=1\), deriva dalla compensazione fra \(C_i\) e il peso del livello;
	non introduce una nuova simmetria gruppale.
	
	\begin{example}[Ricostruzione di un codice]
		Si consideri il codice
		\[
		\boldsymbol{\rho}
		=
		(\rho_2,\rho_1,\rho_0),
		\]
		con
		\[
		\rho_0=(0\,1\,2),
		\qquad
		\rho_1=(0\,2),
		\qquad
		\rho_2=(0\,1),
		\]
		dove i cicli agiscono su \(\{0,1,2\}\). Le statistiche di fibra sono
		\[
		\begin{array}{c@{\qquad}ccc}
			\toprule
			i & \Sigma_{i,0} & \Sigma_{i,1} & \Sigma_{i,2}\\
			\midrule
			0 & 117 & 126 & 108\\
			1 & 144 & 117 &  90\\
			2 & 117 &  36 & 198\\
			\bottomrule
		\end{array}
		\]
		La formula di ricostruzione restituisce le terne dei valori
		\[
		\begin{aligned}
			(\rho_0(0),\rho_0(1),\rho_0(2))
			&=
			\left(
			\frac{117-108}{9},
			\frac{126-108}{9},
			\frac{108-108}{9}
			\right)
			=(1,2,0),\\
			(\rho_1(0),\rho_1(1),\rho_1(2))
			&=
			\left(
			\frac{144-90}{27},
			\frac{117-90}{27},
			\frac{90-90}{27}
			\right)
			=(2,1,0),\\
			(\rho_2(0),\rho_2(1),\rho_2(2))
			&=
			\left(
			\frac{117-36}{81},
			\frac{36-36}{81},
			\frac{198-36}{81}
			\right)
			=(1,0,2).
		\end{aligned}
		\]
		Le tre permutazioni iniziali vengono quindi ricostruite integralmente dai
		nove dati aggregati.
	\end{example}
	
	Nel linguaggio del gioco, il codice di raccolta determina la procedura
	ordinaria e la trasformazione globale; il tabellone ne rappresenta l'azione
	sulle ventisette posizioni. Questa distinzione evita di identificare
	l'esecuzione fisica con la permutazione finale pur conservando, nel modello
	qui considerato, la corrispondenza biunivoca fra i due livelli.
	
	\section{Conseguenze strutturali}
	
	La fattorizzazione separata rende immediatamente leggibili alcuni invarianti
	della trasformazione globale.
	
	\begin{proposition}
		Per
		\[
		\boldsymbol{\rho}
		=(\rho_{m-1},\dots,\rho_0)\in S_b^m
		\]
		valgono
		\[
		\ord(\pi_{\boldsymbol{\rho}})
		=
		\mcm\bigl(
		\ord(\rho_{m-1}),\dots,\ord(\rho_0)
		\bigr)
		\]
		e
		\[
		\fix(\pi_{\boldsymbol{\rho}})
		=
		\prod_{i=0}^{m-1}\fix(\rho_i).
		\]
	\end{proposition}
	
	\begin{proof}
		La potenza \(\pi_{\boldsymbol{\rho}}^n\) è l'identità esattamente quando
		\(\rho_i^n\) è l'identità per ogni \(i\); il minimo intero positivo con
		questa proprietà è il minimo comune multiplo degli ordini locali. Una
		posizione è fissa quando ciascuna delle sue cifre è fissata dalla
		permutazione del livello corrispondente. Le scelte delle cifre sono
		indipendenti, e le cardinalità si moltiplicano.
	\end{proof}
	
	\begin{remark}
		Le classi di coniugio del gruppo interno \(S_b^m\) sono scelte
		indipendentemente nei fattori e sono quindi \(p(b)^m\), dove \(p(b)\) è il
		numero delle partizioni di \(b\). Per \(b=m=3\) se ne ottengono \(3^3=27\).
		Questa classificazione riguarda il gruppo digitalmente separabile; la
		coniugazione nel gruppo ambiente \(S_{b^m}\) può identificare elementi che
		restano distinti in \(S_b^m\).
	\end{remark}
	
	\section{Discussione conclusiva}
	
	Il modello posizionale, la descrizione tensoriale e l'azione di gruppo
	rappresentano tre livelli della stessa struttura. La procedura fisica estrae
	ciclicamente le cifre dell'indirizzo; il codice locale assegna a ciascun
	livello una permutazione; il ciclo completo produce una permutazione globale
	che conserva tutte le partizioni digitali e possiede una fattorizzazione di
	Kronecker unica. La catena può essere riassunta come
	\[
	\begin{aligned}
		\text{procedura fisica}
		&\longrightarrow
		\text{codice locale}
		\longrightarrow
		\text{permutazione digitalmente separabile}
		\\
		&\longrightarrow
		\text{fattorizzazione e ricostruzione}.
	\end{aligned}
	\]
	
	La caratterizzazione mediante le partizioni \(\mathcal F_i\) individua la
	classe globale senza richiamare il gesto concreto del gioco. La formula delle
	statistiche di fibra aggiunge un secondo livello informativo: entro questa
	classe, somme aggregate delle immagini determinano tutti i fattori locali.
	La ricostruzione non dipende quindi dalla conoscenza della storia fisica del
	mazzo, ma dalla permanenza della separabilità digitale nella trasformazione
	finale.
	
	Le ipotesi del modello delimitano con precisione il risultato. L'omogeneità
	\(N=b^m\), la conservazione dell'ordine interno, l'assenza di rovesciamenti e
	la scelta non adattiva delle raccolte rendono possibile l'azione componente
	per componente. L'introduzione di controlli supplementari, dipendenze fra i
	livelli o fattorizzazioni non omogenee conduce a gruppi e dinamiche differenti.
	Il caso rettangolare discusso nell'appendice conserva la permutazione dei
	fattori tensoriali, ma perde la chiusura ciclica di livelli tutti uguali; nel
	gioco delle ventuno carte la determinazione della carta scelta emerge inoltre
	da una contrazione dello spazio delle possibilità, distinta dalla
	permutazione fisica del mazzo.
	
	Il caso \(27=3^3\) non è pertanto speciale soltanto per la sua cardinalità.
	La compatibilità fra scrittura ternaria, rotazione ciclica dei fattori e
	azione separata delle raccolte trasforma un procedimento ricreativo in una
	famiglia rigidamente riconoscibile di permutazioni globali.
	
	\appendix
	
	\section{Il caso rettangolare 
		\texorpdfstring{\(N=i\,j\)}{N=ij} e la dinamica informativa}
	
	La fattorizzazione \(N=b^m\) utilizza livelli omogenei. Nei procedimenti
	rettangolari il numero delle carte assume invece la forma
	\[
	N=i\,j,
	\]
	con fattori che possono essere differenti. Se il mazzo viene distribuito
	ciclicamente in \(i\) colonne di lunghezza \(j\), ogni posizione si scrive
	\[
	x=iq+r,
	\qquad
	0\le r<i,
	\qquad
	0\le q<j.
	\]
	Il resto \(r\) identifica la colonna e il quoziente \(q\) la quota al suo
	interno. Alla posizione si associa
	\[
	e_q^{(j)}\otimes e_r^{(i)}
	\in
	\mathbb R^j\otimes\mathbb R^i.
	\]
	
	Lo scambio fra quota e colonna è realizzato dalla matrice di commutazione
	\[
	K_{j,i}:
	\mathbb R^j\otimes\mathbb R^i
	\longrightarrow
	\mathbb R^i\otimes\mathbb R^j,
	\]
	definita da
	\[
	K_{j,i}
	\bigl(e_q^{(j)}\otimes e_r^{(i)}\bigr)
	=
	e_r^{(i)}\otimes e_q^{(j)}.
	\]
	Vale
	\[
	K_{j,i}^{-1}=K_{i,j}=K_{j,i}^{\mathsf T}.
	\]
	
	Sia \(\tau\in S_i\) la permutazione che assegna alla colonna \(r\) la
	posizione di blocco \(\tau(r)\), e sia \(P_\tau\) la sua matrice. Conservando
	l'ordine interno delle colonne, il passaggio fisico è
	\[
	\pi_\tau(iq+r)
	=
	j\tau(r)+q
	\]
	e il corrispondente operatore è
	\[
	M_\tau
	=
	(P_\tau\otimes I_j)\circ K_{j,i}:
	\mathbb R^j\otimes\mathbb R^i
	\longrightarrow
	\mathbb R^i\otimes\mathbb R^j.
	\]
	Dopo aver identificato i due prodotti tensoriali ordinati con
	\(\mathbb R^{ij}\) mediante le rispettive basi linearizzate, \(M_\tau\) è
	rappresentato da una matrice di permutazione \(ij\times ij\). La sua
	invertibilità descrive la reversibilità del passaggio fisico.
	
	L'operatore \(K_{j,i}\) e la rotazione \(\MSC_m\) esprimono entrambi una
	permutazione dei fattori. Nel caso \(b^m\), tuttavia, tutti i fattori hanno la
	stessa dimensione e una rotazione di ordine \(m\) espone successivamente ogni
	cifra della medesima base. Nel caso rettangolare, quota e colonna hanno in
	generale dimensioni differenti; lo scambio è invertibile, ma non produce da
	solo una chiusura digitale analoga.
	
	Nel gioco delle ventuno carte l'ordine di raccolta è inoltre adattivo. Dopo
	ogni distribuzione il partecipante indica la colonna contenente la carta
	scelta e quella colonna viene collocata al centro. Se, più in generale, la
	colonna indicata viene collocata nella posizione di blocco
	\(c\in\{0,\dots,i-1\}\), la posizione della sola carta scelta segue la mappa
	\[
	F_c(x)
	=
	jc+\left\lfloor\frac{x}{i}\right\rfloor.
	\]
	Questa non è una singola permutazione fisica del mazzo: per ogni posizione
	ipotetica seleziona il ramo corrispondente alla colonna che verrebbe indicata,
	e soddisfa
	\[
	F_c(iq)=F_c(iq+1)=\cdots=F_c(iq+i-1)=jc+q.
	\]
	La mappa è quindi molti-a-uno.
	
	Per il gioco classico delle ventuno carte,
	\[
	i=3,
	\qquad
	j=7,
	\qquad
	c=1,
	\]
	si ha
	\[
	F(x)=7+\left\lfloor\frac{x}{3}\right\rfloor.
	\]
	Partendo da
	\[
	X_0=\{0,1,\dots,20\},
	\]
	l'evoluzione degli insiemi di posizioni compatibili è
	\[
	X_1=\{7,8,\dots,13\},
	\qquad
	X_2=\{9,10,11\},
	\qquad
	X_3=\{10\},
	\]
	e quindi
	\[
	21\longmapsto7\longmapsto3\longmapsto1.
	\]
	La permutazione fisica del mazzo resta invertibile a ogni stadio; ciò che si
	contrae è l'incertezza relativa alla carta scelta. Questa distinzione separa
	la dinamica del mazzo dalla dinamica informativa e segna il confine naturale
	della struttura digitalmente separabile sviluppata nel corpo dell'articolo.
	
	\begin{thebibliography}{9}
		
		\bibitem{bolker}
		E.~D. Bolker,
		\emph{Gergonne's Card Trick, Positional Notation, and Radix Sort},
		Mathematics Magazine \textbf{83} (2010), no.~1, 46--49.
		\href{https://doi.org/10.4169/002557010X479983}{doi:10.4169/002557010X479983}.
		
		\bibitem{champanerkar-jani}
		J. Champanerkar and M. Jani,
		\emph{Using a Card Trick to Illustrate Fixed Points and Stability},
		PRIMUS \textbf{25} (2015), no.~5, 473--483.
		\href{https://doi.org/10.1080/10511970.2015.1031308}{doi:10.1080/10511970.2015.1031308}.
		
		\bibitem{diaconis-graham-kantor}
		P. Diaconis, R.~L. Graham and W.~M. Kantor,
		\emph{The Mathematics of Perfect Shuffles},
		Advances in Applied Mathematics \textbf{4} (1983), 175--196.
		\href{https://doi.org/10.1016/0196-8858(83)90009-X}{doi:10.1016/0196-8858(83)90009-X}.
		
		\bibitem{henderson-searle}
		H.~V. Henderson and S.~R. Searle,
		\emph{The Vec-Permutation Matrix, the Vec Operator and Kronecker Products: A Review},
		Linear and Multilinear Algebra \textbf{9} (1981), no.~4, 271--288.
		\href{https://doi.org/10.1080/03081088108817379}{doi:10.1080/03081088108817379}.
		
		\bibitem{quintero}
		R. Quintero,
		\emph{On a Mathematical Model for an Old Card Trick},
		Recreational Mathematics Magazine \textbf{4} (2017), no.~7, 65--77.
		\href{https://doi.org/10.1515/rmm-2017-0014}{doi:10.1515/rmm-2017-0014}.
		
	\end{thebibliography}
	
\end{document}
